\section{VFS}
\label{sec:vfs}

VFS (\textit{virtual memory system}) 顾名思义，就应该用 \cpp{virtual} 函数实现。

某种程度上来说算是个双关，不过倒也不是开玩笑。Linux 里用的是一大堆函数指针，虚函数用的 vtable 也同样是一大堆函数指针，两者不应该有什么区别。好消息是，vtable 是不需要 C++ runtime 的，所以我的确就这么做了。

不过如果想把父类转为子类的话， \cpp{dynamic_cast} 需要 C++ runtime：它需要启用 -frtti。所幸对于单继承而言，我们有绕过它的方法，这正是 LLVM 所做的。具体的介绍请见 \cref{sec:rtti}。当然，用一个额外的 \cpp{void* data} 确实可以解决一切，但是这也太麻烦了。

因此，我的 inode 基类如下：
\begin{minted}{cpp}
class inode {
  atomic<unsigned> refcnt;
protected:
  atomic<unsigned> lnkcnt;
public:
  virtual ~inode() = default;
  virtual size_t read(size_t offset, void* buf, size_t len, int flags) = 0;
  virtual size_t write(size_t offset, const void* buf, size_t len, int flags) = 0;
  // Other virtual member functions
};
\end{minted}

而每种文件系统的 inode 都会实现对应的 read/write 方法。

可以注意到，我维护了两个数据：refcnt 与 lnkcnt。它们的区别是，refcnt 是这个 inode 被内存中其他结构持有的数量，也就是通常的 shared\_ptr 所维护的那个数；而 lnkcnt 则是其他 inode 连接到它的数量，或者说是文件系统中 inode 的 symlink 数。

同样地，我们有两对函数 drop/ref 以及 unlinked/linked 来更改计数器并且释放资源。

当 refcnt 归零时，只有在这个文件在硬盘上有备份的时候，我们才可以直接在内存中删除它。这是文件系统的属性所决定的，比如 ext2 就认为是有备份的，而 tmpfs 则没有。（我们暂时不会 swap。）如果删除了没有硬盘备份的 inode，之后再想打开的时候，数据就没有了。所以，对这种 inode 而言，只有确认了“无法再次打开”之后才可以删除，而这正是 lnkcnt 所做的事情。

同理，如果 lnkcnt 归零但是有 refcnt，我们也暂时不能删除：可能还有进程在使用它，得等使用完了再删掉。在使用完成后 refcnt 会自然减少，所以不必担心会漏掉。

这样的坏处是，对于 ext2 而言，我们需要维护两份 lnkcnt。一份作为 metadata 写进硬盘，另一份留在内存里，不过只要在创建和删除时注意一下就可以了。
